\section{Introduction}
Continual learning asks how a model can learn a sequence of tasks without forgetting the previous ones. Humans do this naturally, but neural networks trained sequentially suffer from \emph{catastrophic forgetting}: either they fail to learn the new task, or they learn it and lose their performance on the earlier ones. Elastic Weight Consolidation (EWC)~\cite{Kirkpatrick_2017} addresses this trade-off with a quadratic penalty, weighted by the Fisher information matrix, that keeps the parameters close to their values after the previous task. Despite its empirical success, EWC lacks theoretical guarantees. The goal of this project is to study EWC-type methods in an analytically tractable toy model, the multi-task spiked Wigner model, where the performance of spectral estimators and their phase transitions can be computed exactly in the high-dimensional limit. Last semester, we tested EWC on real data, introduced the multi-task spiked model, studied a naive spectral method numerically, and derived two EWC-inspired estimators. This report gives the corresponding theory: we derive the asymptotic overlaps of the naive method and of the two EWC-inspired methods with the signals, and use them to tune and compare the methods.
%\subsection{Summary of the project - goal}
%Study the naïve method, method A and method B.
